% problem-type: nhiet-doan-huu-han-tach-bien
% order: 65
% Nguon: De thi MAT2306 (PTDHR 1), cuoi ky 2023-2024, K66, Cau 4. (co dap an chinh thuc)

\section{Nhiệt trên đoạn hữu hạn --- tách biến chuỗi hàm riêng}

\subsection*{Đề bài ví dụ}

Xét bài toán biên --- ban đầu cho phương trình
\[
u_t(x,t) = 2\,u_{xx}(x,t) + f(x,t), \qquad 0 < x < 1,\ t > 0,
\]
với điều kiện biên
\[
u(0,t) = \mu(t), \qquad u_x(1,t) = \nu(t), \qquad t \ge 0,
\]
và điều kiện ban đầu $u(x,0) = g(x)$, $0 \le x \le 1$.

\begin{enumerate}
\item[(a)] Dùng tích phân năng lượng dạng $I(t) = \int_0^1 u^2(x,t)\,dx$ để chứng minh bài toán đã cho có duy nhất nghiệm.
\item[(b)] Với $\mu(t) = 100$, $\nu(t) = 0$, tìm hàm $v(x)$ thỏa mãn $v''(x) = 0$ và $v(0) = 100$, $v'(1) = 0$. Khi đó hàm $w(x,t) = u(x,t) - v(x)$ thỏa mãn bài toán nào?
\item[(c)] Giải bài toán cho hàm $w(x,t)$ ở câu (b) khi $f(x,t) = e^{-t}$, $g(x) = 100$. Từ đó giải bài toán ban đầu với các dữ liệu $f, g, \mu, \nu$ cho ở câu (b) và (c).
\end{enumerate}

\emph{\small(Nguồn: Đề thi MAT2306 --- PTĐHR 1, cuối kỳ 2023--2024, Câu 4. Có đáp án chính thức.)}

\subsection*{Nhận ra dạng này khi nào?}

\begin{itemize}
\item Phương trình nhiệt (truyền nhiệt/khuếch tán) trên một \textbf{đoạn/thanh hữu hạn} $[0,L]$: $u_t = a\,u_{xx} + f$.
\item Có điều kiện biên ở \textbf{hai đầu} $x=0$ và $x=L$: kiểu Dirichlet ($u$ cho trước) hoặc Neumann ($u_x$ cho trước), có thể trộn lẫn.
\item Biên có thể \textbf{KHÔNG thuần nhất} (vế phải $\ne 0$) hoặc phương trình có \textbf{nguồn} $f \ne 0$.
\item Thường yêu cầu: (a) chứng minh duy nhất nghiệm bằng năng lượng; (b) đưa biên về thuần nhất; (c) giải bằng tách biến.
\end{itemize}

\begin{meobox}
Biên KHÔNG thuần nhất $\Rightarrow$ trừ đi \textbf{nghiệm dừng} $v(x)$ để hàm mới $w=u-v$ có biên thuần nhất. Hàm riêng tùy biên:
\begin{itemize}
\item Hai đầu Dirichlet $\Rightarrow$ $X_n = \sin(n\pi x/L)$.
\item Hai đầu Neumann (cách nhiệt) $\Rightarrow$ $X_n = \cos(n\pi x/L)$.
\item Hỗn hợp Dirichlet--Neumann $\Rightarrow$ $X_n = \sin\!\big((n+\tfrac12)\pi x/L\big)$.
\end{itemize}
\end{meobox}

\subsection*{Công thức \& phương pháp}

\begin{dieukienbox}
Dùng phương pháp này khi: bài nhiệt $u_t = a\,u_{xx} + f$ trên đoạn hữu hạn $[0,L]$, biên Dirichlet/Neumann (có thể không thuần nhất), cần chứng minh duy nhất và/hoặc giải tường minh.
\end{dieukienbox}

\begin{nhobox}
Quy trình 3 bước:
\begin{enumerate}
\item \textbf{Năng lượng $\Rightarrow$ duy nhất.} Cho hiệu $w=u_1-u_2$. Đặt $I(t)=\int_0^L w^2\,dx$. Chỉ cần chứng minh $I$ thoả:
\begin{itemize}
  \item $I(t)\ge 0$
  \item $I(0)=0$
  \item $I'(t)\le 0$
\end{itemize}
$\Rightarrow I\equiv 0 \Rightarrow w\equiv 0 \Rightarrow u_1=u_2$.

\emph{Vì sao có ba điều đó:} $I\ge0$ do bình phương; $I(0)=0$ do dữ kiện đầu $w(\cdot,0)=0$; $I'(t)\le 0$ do $I'=2a\int w\,w_{xx}=-2a\int(w_x)^2\le 0$ (tích phân từng phần, số hạng biên triệt vì $w$ có biên thuần nhất).
\item \textbf{Nghiệm dừng $\Rightarrow$ biên thuần nhất:} giải $v''=0$ (hoặc $a v''+f_\infty=0$) khớp biên KHÔNG thuần nhất. Đặt $w=u-v$ thì $w$ có biên THUẦN NHẤT.
\item \textbf{Tách biến chuỗi hàm riêng:} viết $w=\sum_n T_n(t)X_n(x)$, $X_n$ là hàm riêng Sturm--Liouville của biên thuần nhất. Giải ODE $T_n'+a\lambda_n T_n = (\text{hệ số nguồn})$; hệ số đầu lấy từ khai triển $g-v$.
\end{enumerate}
\end{nhobox}

\subsection*{Đề bài \& bài giải}

\paragraph{(a) Duy nhất nghiệm bằng phương pháp năng lượng.}\leavevmode

\emph{Bước 1 --- dựng hàm hiệu.} Giả sử bài toán có hai nghiệm $u_1, u_2$ cùng dữ kiện. Đặt $w=u_1-u_2$. Do PDE và mọi điều kiện đều \emph{tuyến tính}, hiệu $w$ thoả bài toán \textbf{thuần nhất hoàn toàn} (vế phải, biên, dữ kiện đầu đều $=0$):
\[
w_t = 2\,w_{xx},\qquad 0<x<1,\ t>0,
\]
\[
w(0,t)=0,\quad w_x(1,t)=0,\quad w(x,0)=0.
\]
Mục tiêu: chứng minh $w\equiv0$, tức $u_1=u_2$.

\emph{Bước 2 --- định nghĩa ``năng lượng''.} Gọi
\[
I(t):=\int_0^1 w^2(x,t)\,dx\ \ge 0
\]
là \textbf{tích phân năng lượng} (bình phương chuẩn $L^2$ của $w$ tại thời điểm $t$). Vì $w(x,0)=0$ nên $I(0)=0$. Ta sẽ chỉ ra $I$ \emph{không tăng}, kết hợp $I\ge0$ buộc $I\equiv0$.

\emph{Bước 3 --- tính $I'(t)$.} Đạo hàm dưới dấu tích phân rồi thay $w_t=2w_{xx}$:
\[
I'(t)=\int_0^1 2w\,w_t\,dx=4\int_0^1 w\,w_{xx}\,dx.
\]
Tích phân từng phần ($u=w$, $dv=w_{xx}dx$, suy ra $du=w_x dx$, $v=w_x$):
\[
\int_0^1 w\,w_{xx}\,dx = \big[w\,w_x\big]_{x=0}^{x=1} - \int_0^1 w_x^2\,dx.
\]
Áp biên: $w(0,t)=0$ làm tắt số hạng biên tại $x=0$; $w_x(1,t)=0$ làm tắt số hạng biên tại $x=1$. Vậy phần biên triệt tiêu và
\[
I'(t)=-4\int_0^1 w_x^2\,dx\ \le\ 0,\qquad t>0.
\]

\emph{Bước 4 --- kết luận.} $I$ không âm, không tăng, lại bắt đầu từ $I(0)=0$, do đó $I(t)\equiv0$. Tích phân của một hàm liên tục không âm bằng $0$ kéo theo bản thân hàm $\equiv0$, tức $w(x,t)\equiv0$. Vậy $u_1\equiv u_2$ --- bài toán có duy nhất nghiệm.

\paragraph{(b) Hạ biên không thuần nhất bằng nghiệm dừng.}\leavevmode

\emph{Bước 1 --- ý tưởng.} Biên $u(0,t)=100\ne0$ ngăn ta dùng tách biến trực tiếp (chuỗi hàm riêng đòi biên $=0$). Cách chuẩn: tìm hàm \textbf{nghiệm dừng} $v(x)$ --- tức nghiệm \emph{không phụ thuộc $t$} của phần thuần nhất của PDE --- khớp đúng các biên không thuần nhất; sau đó trừ đi.

\emph{Bước 2 --- giải $v$.} Vì $v$ độc lập $t$ và $f$ không tham gia ở đây, $v$ thoả
\[
v''(x)=0,\qquad v(0)=100,\quad v'(1)=0.
\]
Tích phân hai lần: $v(x)=ax+b$. Áp $v(0)=100$ cho $b=100$; áp $v'(1)=a=0$. Vậy
\[
\boxed{v(x)\equiv 100.}
\]

\emph{Bước 3 --- bài toán cho $w=u-v$.} Đặt $w(x,t):=u(x,t)-v(x)=u(x,t)-100$. Vì $v$ là hằng nên $w_t=u_t$, $w_{xx}=u_{xx}$. Thay vào PDE và các biên:
\[
w_t=2w_{xx}+f(x,t),\qquad 0<x<1,\ t>0,
\]
\[
w(0,t)=u(0,t)-100=0,\quad w_x(1,t)=u_x(1,t)-0=0,
\]
\[
w(x,0)=g(x)-100.
\]
Biên của $w$ đã \textbf{thuần nhất} (cả hai $=0$); chỉ còn dữ kiện đầu và nguồn $f$.

\paragraph{(c) Tách biến --- giải $w$ rồi ráp lại.}\leavevmode

Áp dữ kiện cụ thể: $f(x,t)=e^{-t}$, $g(x)=100$, nên $w(x,0)=g(x)-100=0$.

\emph{Bước 1 --- tìm hàm riêng (bài toán Sturm--Liouville).} Tách biến \emph{thử} cho phần thuần nhất $w_t=2w_{xx}$: viết $w=X(x)T(t)$, thay vào và chia hai vế cho $2XT$:
\[
\frac{T'(t)}{2T(t)}=\frac{X''(x)}{X(x)}=-\lambda\quad(\text{hằng số tách}).
\]
Vế phụ thuộc $x$ kèm biên thuần nhất sinh \textbf{bài toán trị riêng}:
\[
X''(x)+\lambda X(x)=0,\qquad X(0)=0,\quad X'(1)=0.
\]
Một \emph{trị riêng} là số $\lambda$ mà ODE trên có nghiệm $X\not\equiv0$ thoả cả hai biên; nghiệm đó gọi là \emph{hàm riêng}.

\emph{Tìm $\lambda$.} Với $\lambda>0$, đặt $\lambda=\mu^2$, nghiệm tổng quát $X=A\cos(\mu x)+B\sin(\mu x)$. Biên $X(0)=0$ cho $A=0$, còn $X=B\sin(\mu x)$. Biên $X'(1)=B\mu\cos\mu=0$ với $B\ne0$ buộc $\cos\mu=0$, tức
\[
\mu_n=\big(n+\tfrac12\big)\pi,\qquad \lambda_n=\big(n+\tfrac12\big)^2\pi^2,\qquad n=0,1,2,\dots
\]
(Hai trường hợp $\lambda\le0$ chỉ cho $X\equiv0$, loại.) Hàm riêng tương ứng:
\[
X_n(x)=\sin\!\big((n+\tfrac12)\pi x\big).
\]
Hệ $\{X_n\}$ \textbf{trực giao} trên $[0,1]$:
\[
\int_0^1 X_n(x)X_m(x)\,dx=\tfrac12\,\delta_{nm}.
\]

\emph{Bước 2 --- khai triển nghiệm theo chuỗi Fourier hàm riêng.} Tìm $w$ dưới dạng
\[
w(x,t)=\sum_{n=0}^{\infty} T_n(t)\,X_n(x)=\sum_{n=0}^{\infty} T_n(t)\sin\!\big((n+\tfrac12)\pi x\big),
\]
với các \emph{hệ số khai triển} $T_n(t)$ cần tìm (mỗi $T_n$ là một hàm thời gian).

\emph{Bước 3 --- khai triển nguồn $f$ theo cùng hệ hàm riêng.} Viết $f(x,t)=\sum_n f_n(t)X_n(x)$, hệ số (dùng trực giao và $\int_0^1 X_n^2=\tfrac12$):
\[
f_n(t)=\frac{\int_0^1 f(x,t)\sin((n+\tfrac12)\pi x)\,dx}{\int_0^1 \sin^2((n+\tfrac12)\pi x)\,dx}
=2e^{-t}\int_0^1 \sin\!\big((n+\tfrac12)\pi x\big)\,dx.
\]
Tính tích phân trong:
\[
\int_0^1\!\sin\!\big((n+\tfrac12)\pi x\big)\,dx
=\Big[-\tfrac{\cos((n+\tfrac12)\pi x)}{(n+\tfrac12)\pi}\Big]_0^1
=\frac{1-\cos((n+\tfrac12)\pi)}{(n+\tfrac12)\pi}=\frac{1}{(n+\tfrac12)\pi}=\frac{2}{(2n+1)\pi},
\]
vì $\cos((n+\tfrac12)\pi)=0$. Vậy
\[
f_n(t)=\frac{4\,e^{-t}}{(2n+1)\pi}.
\]

\emph{Bước 4 --- ODE cho $T_n(t)$.} Thay khai triển $w=\sum T_n X_n$ vào $w_t=2w_{xx}+f$ và dùng $X_n''=-\lambda_n X_n$:
\[
\sum_n T_n'(t)X_n(x)=\sum_n\big(-2\lambda_n T_n(t)+f_n(t)\big)X_n(x).
\]
Trực giao của $\{X_n\}$ cho phép \emph{đồng nhất từng mode}:
\[
T_n'(t)+2\lambda_n T_n(t)=f_n(t),\qquad \lambda_n=\big(n+\tfrac12\big)^2\pi^2=\tfrac{(2n+1)^2\pi^2}{4}.
\]
Dữ kiện đầu $w(x,0)=0$ cho $T_n(0)=0$ với mọi $n$.

\emph{Bước 5 --- giải ODE tuyến tính bậc một.} Đặt $\alpha_n:=2\lambda_n=\tfrac{(2n+1)^2\pi^2}{2}$. Nhân thừa số tích phân $e^{\alpha_n t}$:
\[
\big(e^{\alpha_n t}T_n\big)'=\frac{4\,e^{(\alpha_n-1)t}}{(2n+1)\pi}.
\]
Tích phân từ $0$ tới $t$ (dùng $T_n(0)=0$):
\[
e^{\alpha_n t}T_n(t)=\frac{4}{(2n+1)\pi}\cdot\frac{e^{(\alpha_n-1)t}-1}{\alpha_n-1}.
\]
Chia hai vế cho $e^{\alpha_n t}$:
\[
T_n(t)=\frac{4}{(2n+1)\pi(\alpha_n-1)}\big(e^{-t}-e^{-\alpha_n t}\big).
\]
Thay $\alpha_n-1=\tfrac{(2n+1)^2\pi^2-2}{2}$, nên $\dfrac{4}{(2n+1)\pi(\alpha_n-1)}=\dfrac{8}{(2n+1)\pi\big[(2n+1)^2\pi^2-2\big]}$:
\[
T_n(t)=\frac{8}{(2n+1)\pi\big[(2n+1)^2\pi^2-2\big]}\Big(e^{-t}-e^{-\frac{(2n+1)^2\pi^2 t}{2}}\Big).
\]

\emph{Bước 6 --- ráp lại $u=v+w$.}

\begin{nhobox}
Nghiệm cuối cùng:
\[
u(x,t)=100+\sum_{n=0}^{\infty}
\frac{8}{(2n+1)\pi\big[(2n+1)^2\pi^2-2\big]}
\Big(e^{-t}-e^{-\frac{(2n+1)^2\pi^2 t}{2}}\Big)
\sin\!\big((n+\tfrac12)\pi x\big).
\]
\end{nhobox}

\begin{luuybox}
Hai biến thể thực tế cùng họ:
\begin{itemize}
\item \textbf{Khử hạng tử phản ứng $-cu$:} nếu phương trình có thêm số hạng $-cu$, đổi biến $v = e^{\beta t}u$ (chọn $\beta$ thích hợp) để khử nó trước khi tách biến. Xem Đề MAT2306 cuối kỳ 2022--2023 Câu 4 và 2021--2022 Câu 4.
\item \textbf{Hai đầu cách nhiệt (Neumann):} thanh/hình chữ nhật hai đầu cách nhiệt cho chuỗi cosine; khi nguồn trùng một mode riêng thì xảy ra cộng hưởng, sinh số hạng dạng $t\,e^{-\lambda t}$. Xem Đề GK MAT3365 (Toán tin) Câu 3.
\end{itemize}
\end{luuybox}
